# Inventario de figuras y tablas detectadas en tfgtfmthesisuam.pdf

## Figuras
- Figura 5.1: Ranking de val_acc final en CIFAR-10 y SVHN.
- Figura 5.2: Análisis de Data Augmentation en CIFAR-10: heatmaps train/val, mejor val_acc y gap.
- Figura 5.3: Dispersión de val_acc frente a unique_pctg en fc1 para CIFAR-10.
- Figura 5.4: Dispersión de val_acc frente al ratio de dispersión en fc1 para CIFAR-10.
- Figura 5.5: Dispersión de val_acc frente a unique_pctg en fc1 para SVHN.
- Figura 5.6: Dispersión de val_acc frente a unique_pctg en fc1 para FashionMNIST.
- Figura 5.7: Evolución temporal de unique_pctg en conv1, conv3 y fc1.
- Figura 5.8: Curvas de degradación ante ruido gaussiano en datos de entrada.
- Figura 5.9: Curvas de degradación ante ruido en pesos de conv1 y conv3.
- Figura 5.10: Curvas de degradación ante ruido en pesos de fc1.
- Figura 5.11: Retención de accuracy a sigma=0.15 por método y capa.
- Figura 5.12: Radar multidimensional de accuracy, robustez a datos y robustez a pesos.

## Tablas
- Tabla 4.1: Arquitectura CNN utilizada en los experimentos (metodológica, no ampliada por CSV).
- Tabla 5.1: unique_pctg y entropía H en fc1 comparando bins disponibles.
- Tabla 5.2: Tres métricas de activación en fc1 para el escenario óptimo.
- Tabla 5.3: Retención ante ruido máximo en datos de entrada.
- Tabla 5.4: Retención a sigma=0.15 de ruido en pesos por capa.
- Tabla 5.5: Mejor y peor método por dimensión de evaluación.
- Tabla 5.6: Comparativa multidimensional de regularizadores.

Nota: las tablas 5.1-5.6 se generan extendidas por dataset. Cuando un CSV no contiene una métrica o resolución, se marca como n/d.